\renewcommand{\backrefpagesname}{Citado na(s) página(s):~}
\renewcommand{\backref}{}
\renewcommand*{\backrefalt}[4]{
	\ifcase #1 %
		Nenhuma citação no texto.%
	\or
		Citado na página #2.%
	\else
		Citado #1 vezes nas páginas #2.%
	\fi}%
% ---


% Marca trechos que ainda precisam ser completados pelo autor
\newcommand{\pendente}[1]{\textcolor{red}{[PENDENTE: #1]}}
% Versão curta para células de tabela; a busca por \pendente também a encontra
\newcommand{\pendentecelula}[1]{\textcolor{red}{[#1]}}

% Identificadores de código (tabelas, colunas, caminhos) em monoespaçado, com
% quebra de linha permitida em _ e /. Uso: \cod{ouro_reloginho_indicadores}.
% Não usar \path, que colide com o TikZ.
% O abntex2cite define \UrlLeft{<} e \UrlRight{>} para URLs; aqui eles são
% esvaziados para que identificadores não saiam entre colchetes angulares.
\DeclareUrlCommand\cod{\urlstyle{tt}\def\UrlLeft{}\def\UrlRight{}}
\makeatletter
\g@addto@macro{\UrlBreaks}{\do\_\do\/\do\-}
\Urlmuskip=0mu plus 2mu
\makeatother
% Folga extra para parágrafos que, sem ela, estourariam a margem
\setlength{\emergencystretch}{2em}

% Floats: páginas só de floats começam no topo, e tabelas médias ficam no
% texto em vez de irem sozinhas para uma página.
\makeatletter
\setlength{\@fptop}{0pt}
\setlength{\@fpsep}{12pt plus 1fil}
\makeatother
\renewcommand{\topfraction}{0.85}
\renewcommand{\bottomfraction}{0.6}
\renewcommand{\textfraction}{0.12}
\renewcommand{\floatpagefraction}{0.75}
\setcounter{topnumber}{3}
\setcounter{totalnumber}{4}
